% Generated by `taskgen course examples` from number_theory.crt_general@solvable (seed 6).
% Do not edit: change examples.yaml and regenerate.
\begin{example}[модули не взаимно просты]\label{ex:crt-general}
Выяснить, совместна ли система сравнений, и если да, найти все её решения (модули не обязательно попарно взаимно просты):\[\left\{\begin{array}{l}x \equiv 24\ (\textrm{mod}\ 36) \\ x \equiv 6\ (\textrm{mod}\ 14)\end{array}\right.\]

\par\noindent\textit{Решение.} Критерий: система $x \equiv r_i \pmod{m_i}$ совместна тогда и только тогда, когда $r_i \equiv r_j \pmod{\gcd(m_i, m_j)}$ для всех $i < j$; решение единственно по модулю $\operatorname{lcm}(m_1, \dots, m_k)$.\[\begin{array}{c|c|r|c}(m_i, m_j) & \gcd & r_i - r_j & \gcd \mid r_i - r_j? \\ \hline (36, 14) & 2 & 18 & \text{да} \\ \end{array}\]Все условия выполнены, система совместна. Решаем последовательно: из первого сравнения $x = 24 + 36t$.\par\noindent Присоединяем $x \equiv 6\ (\textrm{mod}\ 14)$: $36t \equiv 6 - 24 = -18\ (\textrm{mod}\ 14)$, $d = \gcd(36, 14) = 2$; делим на $d$: $18t \equiv -9\ (\textrm{mod}\ 7)$, т.\,е. $4t \equiv 5\ (\textrm{mod}\ 7)$.  Обратный к $4$ по модулю $7$ (расширенный алгоритм Евклида для $(7, 4)$):
\[\begin{array}{r|r|r|r}
r_i & q_i & u_i & v_i \\ \hline
7 &  & 1 & 0 \\
4 &  & 0 & 1 \\
3 & 1 & 1 & -1 \\
1 & 1 & -1 & 2 \\
0 & 3 & 4 & -7 \\
\end{array}\qquad r_i = u_i \cdot 7 + v_i \cdot 4\]

$4^{-1} \equiv 2\ (\textrm{mod}\ 7)$.
 $t \equiv 5 \cdot 2 \equiv 3\ (\textrm{mod}\ 7)$, $x = 24 + 36 \cdot 3 = 132$, новый модуль $\operatorname{lcm}(36, 14) = 36 \cdot 7 = 252$: $x \equiv 132\ (\textrm{mod}\ 252)$.\par Проверка: $132 = 36 \cdot 3 + 24$; $132 = 14 \cdot 9 + 6$.\[x \equiv 132\ (\textrm{mod}\ 252).\]
\end{example}
